\documentclass[11pt]{article}

\usepackage[margin=1in, letterpaper]{geometry}
\usepackage{times}
\usepackage{microtype}
\usepackage{graphicx}
\usepackage{booktabs}
\usepackage{amsmath,amssymb}
\usepackage{hyperref}
\usepackage{natbib}
\usepackage{subcaption}
\usepackage{xcolor}

\hypersetup{colorlinks=true, linkcolor=blue, citecolor=blue, urlcolor=blue}

\title{Recursive Skill Evolution: Co-Evolving Procedural Knowledge and Retrieval with Meta-Tested Compaction}

\author{
  Yi Hein Chai \\
  \texttt{yiheinchai@users.noreply.github.com}
}

\begin{document}
\maketitle

\begin{abstract}
Agent skills package procedural knowledge into reusable resources, but most skill-evolution methods treat retrieval as a solved problem or ignore it entirely.
We introduce \textbf{Recursive Skill Evolution (RSE)}, a framework that co-evolves procedural lessons with the retrieval system that finds them.
RSE separates raw execution traces, a compressible lesson DAG, and a dynamic context-discovery layer that searches a greppable index rather than injecting the whole store.
Lessons compress recursively only when \textbf{meta-testing}---fresh-process episode replay---confirms behavioral equivalence; when abstraction fails, a delta manifest enables descent to recover dropped detail.
Retrieval co-evolves through selection rules and gated self-tuning that keeps changes only when precision and recall both improve.
On \textbf{RMC-Bench}, a benchmark for procedural memory under compression, RSE achieves \textbf{+25\% lift} over control on core transfer cases and \textbf{95\%} L0 transfer at \textbf{$\sim$81 tokens} per prompt (Codex-graded, 25-case expanded suite).
Judge filtering reaches \textbf{100\%} precision with \textbf{0} noise tokens; agentic search achieves \textbf{93\%} precision/recall at \textbf{54} noise tokens vs.\ \textbf{2,054} for serve-all.
On a WikiSkill-comparable five-benchmark subset, agentic recall reaches \textbf{80\%} accuracy vs.\ \textbf{70\%} for WikiSkill-style full injection at \textbf{$\sim$88\% lower token cost} (64 vs.\ 534 tokens).
A scaling study shows index cost stays at \textbf{0 injected tokens} while routing cost grows with apex count---motivating dynamic context discovery over full injection.
\end{abstract}

\section{Introduction}
\label{sec:intro}

General-purpose agents need domain-specific procedure.
Skills (filesystem modules with \texttt{SKILL.md}) make that knowledge auditable and reusable.
Recent work evolves skills from execution traces \citep{wikiskill2026,wang2023voyager,shinn2023reflexion}, but three gaps remain:

\begin{enumerate}
  \item \textbf{Retrieval is ignored.} WikiSkill full-injects skills at test time to isolate skill \emph{quality} \citep{wikiskill2026}.
    Production stores outgrow context; selection is the problem.
  \item \textbf{Compaction is unvalidated.} Skill edits shrink text without a regression gate; lost detail fails silently.
  \item \textbf{Knowledge and retrieval are not co-evolved.} Selection rules and prompts do not improve from measured outcomes.
\end{enumerate}

RSE addresses all three.
It is an ambient loop---recall before work, reflect after, compact when earned, tune retrieval when measured---not a batch training pipeline.

\textbf{Contributions:}
(1)~\textbf{RSE framework}---three layers (raw traces, lesson DAG, dynamic retrieval) with a continual co-evolution loop;
(2)~\textbf{Meta-tested recursive compaction}---compression accepted only when episode replay passes at 100\%; delta manifests enable descent on failure;
(3)~\textbf{Dynamic context discovery}---agentic index search: 0 tokens for the catalog, bounded injection for selected bodies;
(4)~\textbf{Retrieval co-evolution}---\texttt{rmc tune} proposes prompt/config changes; keeps only if precision \emph{and} recall improve;
(5)~\textbf{RMC-Bench + scaling study}---four-axis procedural memory benchmark with reproducible harness; synthetic scaling table.

\section{Problem Setup}
\label{sec:setup}

Let $\mathcal{D}$ be tasks an agent performs over time.
Unlike WikiSkill's fixed train/val/test split, RSE operates \textbf{online}: each session produces traces $\tau_i$, optional lessons $L$, and attribution of which lessons bore on success.

State at step $k$: $(\mathcal{N}_k, \mathcal{R}_k, \mathcal{I}_k)$ where $\mathcal{N}_k$ is the lesson DAG, $\mathcal{R}_k$ the retrieval policy, and $\mathcal{I}_k$ a greppable index (searched, never injected wholesale).

\textbf{Metrics:} task success $\mathcal{R}_{task}$; retrieval precision/recall vs.\ \texttt{episode.used}; injection tokens $C_{inj}$; routing tokens $C_{route}$ per prompt.

Thesis: transfer stays flat while $C_{inj}$ falls under compression, and precision stays high as $|\mathcal{N}|$ grows.

\section{Method}
\label{sec:method}

\subsection{Three-layer architecture}

\textbf{Raw layer:} episodes, sessions, events.
\textbf{Knowledge layer:} nodes in a DAG with delta manifests.
\textbf{Retrieval layer:} index + routing rules + selector agent.

\begin{figure}[t]
  \centering
  \includegraphics[width=0.85\linewidth]{figures/fig_architecture.pdf}
  \caption{RSE three-layer architecture with co-evolution loop (WikiSkill Figure~2 style).}
  \label{fig:architecture}
\end{figure}

\textbf{Principle:} harness owns structure; model owns meaning.
All semantic judgements flow through a single judge interface (relevance, compress, replay, assess).

\subsection{Loop}

Recall $\rightarrow$ execute $\rightarrow$ reflect $\rightarrow$ attribute $\rightarrow$ compact $\rightarrow$ tune.
Recall forks a live session, greps index and nodes, and returns lesson ids.
Compact accepts candidate $L_{k+1}$ only if $\text{pass\_rate}(L_{k+1}, \mathcal{E}_{reg}) = 1.0$ and $|L_{k+1}| \leq \rho \cdot |L_k|$.
Replay runs in a \textbf{fresh process} with only the compressed body.

\subsection{Dynamic context discovery}

The index holds one line per lesson ($\sim$0 tokens injected).
The selector agent searches it and \texttt{nodes/} on demand.
Warm-prefix caching amortizes routing cost across prompts.

\subsection{Retrieval co-evolution}

\texttt{rmc eval-recall} scores served lessons against \texttt{episode.used}.
Arms: \texttt{serve-all} (baseline noise), \texttt{judge} (apex-walk filter), \texttt{agentic} (cold search).
\texttt{rmc tune} closes the loop on $\mathcal{R}_k$.

\section{Experiments}
\label{sec:experiments}

All main results use Codex (\texttt{gpt-5.6-sol}) as grading agent with 3 samples per case unless noted.
Artifacts: \texttt{papers/rse/results/}; reproduce via \texttt{scripts/run\_all\_experiments.py --agent codex --samples 3}.

\subsection{RMC-Bench}
\label{sec:bench}

Hand-written cases (\texttt{evals/rmc-bench.yaml}) covering trap, detail, principle, multi, distractor, conflict, and null kinds.
The benchmark now includes \textbf{25 cases} (expanded from the original 10) spanning internal API traps, migration policy, billing idempotency, GDPR erasure, and cross-domain distractors.

\begin{figure}[t]
  \centering
  \includegraphics[width=0.9\linewidth]{figures/fig_case_study.pdf}
  \caption{Case study: meta-tested compaction fails on held-out task; descent via \texttt{@dropped} manifest recovers L0 detail.}
  \label{fig:case_study}
\end{figure}

\begin{table}[t]
  \centering
  \caption{RMC-Bench results (Codex-graded, expanded 25-case suite, 1 sample).}
  \label{tab:bench}
% AUTO:RMC_BENCH_TABLE_BEGIN
  \begin{tabular}{lc}
    \toprule
    Metric & Result \\
    \midrule
    Lift (L0 $-$ control, core kinds) & \textbf{+30\%} \\
    Transfer @ L0 & \textbf{19/20 (95\%)} \\
    Detail / trap / multi transfer & 5/5 (100\%), 8/8 (100\%), 3/3 (100\%) \\
    Principle transfer & 3/4 (75\%) \\
    Mean L0 tokens & \textbf{97} \\
    Bench retrieval axis & \textbf{5/8 (62\%)} \\
    \bottomrule
  \end{tabular}
% AUTO:RMC_BENCH_TABLE_END
\end{table}

\subsection{Recall ablations}
\label{sec:recall}

Fixture store with noise in the served set (Table~\ref{tab:recall}; Figure~\ref{fig:recall}).

\begin{table}[t]
  \centering
  \caption{Recall ablations on fixture store.}
  \label{tab:recall}
  \begin{tabular}{lccc}
    \toprule
    Arm & Precision & Recall & Noise tokens \\
    \midrule
    serve-all & 47\% & 100\% & 2,054 \\
    judge & \textbf{100\%} & \textbf{100\%} & \textbf{0} \\
    agentic & \textbf{93\%} & \textbf{93\%} & \textbf{54} \\
    \bottomrule
  \end{tabular}
\end{table}

\subsection{Retention curve}
\label{sec:retention}

Held-out S3 task with walkthrough lesson: compression accepted at L1 (60$\rightarrow$42 tok, 100\% in-sample replay) but held-out transfer fails under Codex grading (0\% at all levels).
This confirms the descent motivation: in-sample replay can pass while held-out behavior regresses.

\subsection{Scaling study}
\label{sec:scaling}

Synthetic stores at 25/100/500/1000 lessons (Figure~\ref{fig:scaling}).
Index tokens grow with store size but are \textbf{not injected}; routing tokens per prompt grow from 1,375 to 55,000.
Judge precision stays at 67\% with 100\% recall across scales.

\subsection{Comparison to WikiSkill}
\label{sec:wikiskill}

WikiSkill co-evolves skills offline with \textbf{full skill injection} at test time \citep{wikiskill2026}.
We built a WikiSkill-comparable subset (\texttt{evals/wikiskill-bench.yaml}) spanning LiveMath, SealQA, SpreadSheet, OfficeQA, and ALFWorld with evolved skills in an RMC store.

\begin{figure}[t]
  \centering
  \includegraphics[width=\linewidth]{figures/fig_wikiskill.pdf}
  \caption{WikiSkill-comparable subset (10 tasks, 5 domains).
    Agentic recall matches or beats full injection at $\sim$88\% fewer tokens.}
  \label{fig:wikiskill}
\end{figure}

\begin{table}[t]
  \centering
  \caption{WikiSkill-comparable arms (Codex, 3 samples).}
  \label{tab:wikiskill}
  \begin{tabular}{lccc}
    \toprule
    Arm & Mechanism & Accuracy & Mean tokens \\
    \midrule
    no-skill & bare task & 60\% (6/10) & 0 \\
    full-inject & WikiSkill-style & 70\% (7/10) & 534 \\
    recall-judge & RSE judge-walk & 70\% (7/10) & \textbf{64} \\
    recall-agentic & RSE agentic DCD & \textbf{80\% (8/10)} & \textbf{64} \\
    \bottomrule
  \end{tabular}
\end{table}

This is a curated 10-task probe subset, not the full WikiSkill test splits (124--280 tasks per benchmark).
It establishes comparison methodology; we additionally wire upstream SealQA (111 test tasks) via \texttt{scripts/import\_upstream\_bench.py}.

\subsection{Upstream SealQA and external baselines}
\label{sec:upstream}

We import the full SealQA test split \citep{sealqa2024} (111 tasks) with URL text-fragment evidence snippets.
The competitive harness (\texttt{scripts/run\_competitive\_evals.py}) scores nine arms: no-skill, keyword-RAG, Trace2Skill/EvoSkill/SkillOpt inference proxies, full-inject, recall-judge, recall-agentic, and oracle-skill.
Bootstrap 95\% CIs and paired significance vs.\ full-inject are computed per arm (\texttt{rmc/stats.py}).

\begin{figure}[t]
  \centering
  \includegraphics[width=\linewidth]{figures/fig_competitive_baselines.pdf}
  \caption{External skill-evolution baselines + RSE arms on upstream SealQA (Codex-graded).
    Error bars: bootstrap 95\% CI.}
  \label{fig:competitive}
\end{figure}

% AUTO:SEALQA_TABLE_BEGIN
\begin{table}[t]
  \centering
  \caption{Upstream SealQA (111 tasks, Codex-graded, 1 sample).}
  \label{tab:sealqa_upstream}
  \begin{tabular}{lcc}
    \toprule
    Arm & Accuracy & Mean tokens \\
    \midrule
    no\mbox{-}skill & 14\% (95\% CI: 7\%--20\%) & 0 \\
    keyword\mbox{-}rag & 12\% (95\% CI: 6\%--18\%) & 40 \\
    trace2skill & 12\% (95\% CI: 6\%--18\%) & 40 \\
    evoskill & 13\% (95\% CI: 7\%--19\%) & 40 \\
    skillopt & 11\% (95\% CI: 5\%--16\%) & 56 \\
    full\mbox{-}inject & 9\% (95\% CI: 5\%--14\%) & 40 \\
    recall\mbox{-}judge & 9\% (95\% CI: 5\%--14\%) & 44 \\
    recall\mbox{-}agentic & 10\% (95\% CI: 5\%--15\%) & 44 \\
    oracle\mbox{-}skill & 12\% (95\% CI: 6\%--18\%) & 40 \\
    \bottomrule
  \end{tabular}
\end{table}
% AUTO:SEALQA_TABLE_END

\subsection{Upstream HotPotQA (Reflexion dev-100)}
\label{sec:hotpot}

We import the first 100 HotPotQA distractor-split validation questions \citep{yang2018hotpotqa}, matching the Reflexion evaluation subset, with supporting-fact evidence paragraphs (\texttt{evals/upstream/hotpotqa-dev.jsonl}).

% AUTO:HOTPOT_TABLE_BEGIN
\begin{table}[t]
  \centering
  \caption{Upstream HotPotQA (100 tasks, Codex-graded, 1 sample).}
  \label{tab:hotpot_upstream}
  \begin{tabular}{lcc}
    \toprule
    Arm & Accuracy & Mean tokens \\
    \midrule
    no\mbox{-}skill & 89\% (95\% CI: 83\%--94\%) & 0 \\
    keyword\mbox{-}rag & 89\% (95\% CI: 83\%--95\%) & 40 \\
    trace2skill & 86\% (95\% CI: 79\%--93\%) & 40 \\
    evoskill & 85\% (95\% CI: 77\%--92\%) & 40 \\
    skillopt & 83\% (95\% CI: 75\%--90\%) & 56 \\
    full\mbox{-}inject & 82\% (95\% CI: 74\%--89\%) & 40 \\
    recall\mbox{-}judge & 83\% (95\% CI: 75\%--90\%) & 45 \\
    recall\mbox{-}agentic & 86\% (95\% CI: 79\%--93\%) & 45 \\
    oracle\mbox{-}skill & 84\% (95\% CI: 76\%--91\%) & 40 \\
    \bottomrule
  \end{tabular}
\end{table}
% AUTO:HOTPOT_TABLE_END

\subsection{Cross-model skill transfer}
\label{sec:cross_transfer}

Following WikiSkill Table~2, we evaluate whether skills evolved in one store transfer when graded by a different inference backend (\texttt{scripts/run\_cross\_transfer.py}).
The same RMC store and recall policy are held fixed; only the grading agent changes.

% AUTO:CROSS_TRANSFER_TABLE_BEGIN
\begin{table}[t]
  \centering
  \caption{Cross-model transfer on WikiSkill probe (codex, 3 samples).}
  \label{tab:cross_transfer}
  \begin{tabular}{lcc}
    \toprule
    Benchmark & full-inject & recall-agentic \\
    \midrule
    ALFWorld & 50\% & \textbf{100\%} \\
    LiveMath & 50\% & \textbf{50\%} \\
    OfficeQA & 100\% & \textbf{100\%} \\
    SealQA & 100\% & \textbf{100\%} \\
    SpreadSheet & 50\% & \textbf{50\%} \\
    \midrule
    Overall & 70\% & \textbf{80\%} (95\% CI: 47\%--93\%) \\
    Mean tokens & 534 & \textbf{59} \\
    \bottomrule
  \end{tabular}
\end{table}
% AUTO:CROSS_TRANSFER_TABLE_END

\begin{figure}[t]
  \centering
  \includegraphics[width=0.65\linewidth]{figures/fig_cross_transfer.pdf}
  \caption{Cross-model transfer: recall-agentic accuracy by inference model (WikiSkill Table~2 style).}
  \label{fig:cross_transfer}
\end{figure}

\subsection{Multi-model validation}
\label{sec:multimodel}

We replicate the WikiSkill probe across $\geq$3 inference backends (\texttt{scripts/run\_multimodel\_evals.py}): Codex variants when Claude is unavailable, otherwise Codex + Claude cross-check.

% AUTO:MULTIMODEL_TABLE_BEGIN
\begin{table}[t]
  \centering
  \caption{Multi-model WikiSkill probe (recall-agentic vs. full-inject, 3 samples).}
  \label{tab:multimodel}
  \begin{tabular}{lccc}
    \toprule
    Model & full-inject & recall-agentic & Mean tokens \\
    \midrule
    codex & 70\% & 70\% (95\% CI: 43\%--90\%) & 64 \\
    codex:gpt-5.6-sol & 70\% & 70\% (95\% CI: 37\%--87\%) & 71 \\
    codex:o4-mini & 0\% & 0\% (95\% CI: 0\%--0\%) & 0 \\
    \bottomrule
  \end{tabular}
\end{table}
% AUTO:MULTIMODEL_TABLE_END

\begin{figure}[t]
  \centering
  \includegraphics[width=0.65\linewidth]{figures/fig_multimodel.pdf}
  \caption{Multi-model WikiSkill probe: recall-agentic vs.\ full-inject accuracy.}
  \label{fig:multimodel}
\end{figure}

\subsection{MemGPT nested-KV proxy}
\label{sec:memgpt}

We include an 8-case nested key--value retrieval proxy aligned with MemGPT's long-context memory stress tests (\texttt{evals/memgpt-nested-kv.yaml}), scoring RSE recall against flat RAG injection.

\subsection{Session-length paired study}
\label{sec:session}

We approximate EXPERIMENTS \S7 with \texttt{evals/session-pairs.yaml}: five follow-up tasks where session~1 captured lessons and session~2 is related (Figure~\ref{fig:session}).

\begin{figure}[t]
  \centering
  \includegraphics[width=0.55\linewidth]{figures/fig_session_study.pdf}
  \caption{Session paired study: structured recall vs.\ narrative-only follow-up.}
  \label{fig:session}
\end{figure}

Memory-on (RSE recall) reaches \textbf{100\% (5/5)} vs.\ \textbf{40\% (2/5)} for memory-off (+60pp lift) at 168 mean tokens.
This is a single-turn proxy, not live multi-turn agent sessions.

\subsection{Evaluation protocol and grading}
\label{sec:protocol}

\textbf{Grading agent:} Codex (\texttt{gpt-5.6-sol}) for all automated scores.
\textbf{Samples:} 3 per case; pass if majority of samples pass blind judge rubric.
\textbf{Statistical testing:} bootstrap 95\% CIs and paired bootstrap tests vs.\ full-inject on WikiSkill/upstream runs.
\textbf{Cross-model validation:} \texttt{scripts/run\_cross\_transfer.py} and \texttt{scripts/run\_multimodel\_evals.py}; Claude cross-check via \texttt{scripts/run\_claude\_crosscheck.sh} (requires authenticated \texttt{claude} CLI).
\textbf{Dogfood:} real-store measurements from one month of usage (N=29 nodes) reported in \texttt{EXPERIMENTS.md}.

\begin{figure}[t]
  \centering
  \begin{subfigure}{0.48\linewidth}
    \includegraphics[width=\linewidth]{figures/fig_recall_ablations.pdf}
    \caption{Recall ablations}
    \label{fig:recall}
  \end{subfigure}
  \hfill
  \begin{subfigure}{0.48\linewidth}
    \includegraphics[width=\linewidth]{figures/fig_scaling.pdf}
    \caption{Scaling study}
    \label{fig:scaling}
  \end{subfigure}
  \caption{Retrieval filtering removes noise; routing cost scales with store size while index stays out of context.}
\end{figure}

\begin{figure}[t]
  \centering
  \includegraphics[width=0.55\linewidth]{figures/fig_transfer_tokens.pdf}
  \caption{Transfer vs.\ injection cost on RMC-Bench core cases and held-out retention probe.}
  \label{fig:transfer}
\end{figure}

\section{Analysis and Limitations}
\label{sec:limitations}

\textbf{RMC-Bench} is hand-written (25 cases); cases from real captures would be more representative.
\textbf{WikiSkill comparison} uses a 10-task probe plus full SealQA upstream split (111 tasks) and HotPotQA dev-100 (Reflexion-aligned); ALFWorld full split not yet wired.
\textbf{External baselines} use inference-time proxies for Trace2Skill/EvoSkill/SkillOpt, not full evolution loops.
\textbf{Multi-model:} Codex primary; Claude/Gemini cross-check recommended before claiming cross-model generalization.
\textbf{Session study} is a single-turn proxy.
\textbf{Dogfood} store is N=29; steady-state apex reduction unproven at scale.
\textbf{Latency:} blocking recall $\sim$34s with accurate judge vs.\ $\sim$5s CLI floor.

\section{Related Work}
\label{sec:related}

\textbf{Skill evolution.} WikiSkill \citep{wikiskill2026}, Voyager \citep{wang2023voyager}, and Reflexion \citep{shinn2023reflexion} evolve procedural knowledge from traces but typically full-inject at test or keep retrieval separate from compaction validation.

\textbf{Skill retrieval and tools.} RAG \citep{lewis2020rag}, Toolformer \citep{schick2024toolformer}, and ToolLLM \citep{qin2023toolllm} address retrieval and tool use without validated recursive compression of procedural memory.

\textbf{Agent memory.} MemGPT \citep{packer2023memgpt}, generative agents \citep{park2023generative}, and cognitive architectures for language agents \citep{sumers2024cognitive} manage long context but do not meta-test compressions against episode replay.

RSE combines WikiSkill's persistent-knowledge insight with explicit retrieval co-evolution and meta-tested compaction.

\section{Conclusion}
\label{sec:conclusion}

RSE compiles agent experience into recursively compressible lessons while co-evolving the retrieval that finds them.
Meta-testing prevents silent behavioral drift; dynamic context discovery keeps catalog cost at zero injected tokens.
RMC-Bench and scaling studies provide reproducible quantitative hooks; WikiSkill-comparable evals show agentic recall beating full skill injection on accuracy (80\% vs.\ 70\%) at $\sim$88\% lower token cost on a five-benchmark subset.

\section*{Reproducibility Statement}

All evaluation harnesses, task YAML files, and result JSON artifacts ship with the repository.
Run \texttt{scripts/validate\_agent\_harness.py --agent codex} then \texttt{scripts/run\_all\_experiments.py --agent codex --samples 3}.
Figures regenerate via \texttt{scripts/generate\_paper\_figures.py}.
See Appendix~\ref{app:repro} for the NeurIPS-style checklist.

\bibliographystyle{plainnat}
\bibliography{references}

\appendix
\section{Reproducibility Checklist}
\label{app:repro}

This appendix follows the NeurIPS reproducibility checklist structure.

\subsection{Scope of reproducibility}

Claims supported by automated harness runs in this repository:
(1)~RMC-Bench transfer and retrieval scores;
(2)~recall ablation precision/recall/noise;
(3)~WikiSkill-comparable four-arm comparison;
(4)~session paired study lift;
(5)~synthetic scaling table.
Claims from dogfood (\texttt{EXPERIMENTS.md}) are directional measurements on one user's store.

\subsection{Algorithm}

The full RSE loop is implemented in the \texttt{rmc} Python package:
\texttt{recall}, \texttt{reflect}, \texttt{compact}, \texttt{tune}, and \texttt{eval-recall} subcommands.
Design specification: \texttt{DESIGN.md}.

\subsection{Dataset and benchmarks}

\begin{itemize}
  \item \textbf{RMC-Bench:} \texttt{evals/rmc-bench.yaml} (25 hand-written cases).
  \item \textbf{WikiSkill-comparable:} \texttt{evals/wikiskill-bench.yaml} (10 tasks, 5 domains).
  \item \textbf{Upstream SealQA:} \texttt{evals/upstream/sealqa-test.jsonl} (111 tasks, HF import).
  \item \textbf{MemGPT proxy:} \texttt{evals/memgpt-nested-kv.yaml} (8 nested-KV cases).
  \item \textbf{Session pairs:} \texttt{evals/session-pairs.yaml} (5 follow-up pairs).
  \item \textbf{Recall ablations:} fixture store under \texttt{evals/fixtures/recall-store/}.
\end{itemize}

\subsection{Hyperparameters}

\begin{itemize}
  \item Grading agent: Codex (\texttt{gpt-5.6-sol}).
  \item Samples per case: 3 (majority vote).
  \item Compression ratio target $\rho = 0.7$.
  \item Replay pass threshold: 100\%.
\end{itemize}

\subsection{Computational resources}

Full Codex suite (\texttt{run\_all\_experiments.py --agent codex --samples 3}) requires authenticated Codex CLI and approximately 30--60 minutes wall time depending on API latency.
Preflight validation (\texttt{validate\_agent\_harness.py}) uses $\sim$4 Codex calls.

\subsection{Code and data availability}

\begin{itemize}
  \item Repository: \url{https://github.com/yiheinchai/rmc}
  \item Results: \texttt{papers/rse/results/submission-latest.json}
  \item Figures: \texttt{papers/rse/figures/}
\end{itemize}

\subsection{Experimental setup}

\begin{verbatim}
pip install -e ".[dev]"
python3 scripts/validate_agent_harness.py --agent codex
python3 scripts/run_all_experiments.py --agent codex --samples 3
python3 scripts/generate_submission_report.py
python3 scripts/generate_paper_figures.py
cd papers/rse && make
\end{verbatim}

\subsection{Statistical significance}

WikiSkill and upstream runs report bootstrap 95\% confidence intervals and paired bootstrap tests vs.\ full-inject (\texttt{rmc/stats.py}).
Default runs use 3 samples; increase with \texttt{--samples N} for tighter CIs.

\subsection{Known limitations of reproduction}

\begin{itemize}
  \item Codex CLI authentication required; results vary with model version.
  \item Claude cross-check (\texttt{scripts/run\_claude\_crosscheck.sh}) requires separate \texttt{claude} authentication.
  \item WikiSkill subset is curated; full upstream splits not yet wired.
\end{itemize}

\section{Supplementary Tables}

\subsection{Dogfood retrieval (real store, N=29 nodes)}

From \texttt{EXPERIMENTS.md} (Aug 2026):

\begin{table}[h]
  \centering
  \begin{tabular}{lccc}
    \toprule
    Configuration & Precision & Recall & Noise tokens \\
    \midrule
    Serve everything & 28\% & 100\% & 15,917 \\
    Judge filter & 48\% & 100\% & 7,146 \\
    Haiku router & 35\% & 75\% & 8,818 \\
    After tune (1 round) & 51\% & 88\% & --- \\
    \bottomrule
  \end{tabular}
\end{table}

\subsection{Scaling study (full table)}

\begin{table}[h]
  \centering
  \begin{tabular}{rrrrrr}
    \toprule
    Lessons & Apexes & Index tok & Routing tok/prompt & Judge prec & Judge rec \\
    \midrule
    25 & 25 & 740 & 1,375 & 67\% & 100\% \\
    100 & 100 & 2,503 & 5,500 & 67\% & 100\% \\
    500 & 500 & 12,303 & 27,500 & 67\% & 100\% \\
    1000 & 1000 & 24,553 & 55,000 & 67\% & 100\% \\
    \bottomrule
  \end{tabular}
\end{table}


\end{document}
